\unnumchapter{Постановка задачи}

Целью данной работы является исследование механизмов повышения надёжности
сервиса в федерации Kubernetes-кластеров: разработка сквозного контура
управления, связывающего наблюдаемость, политики размещения, выбор кластера и
аварийное восстановление, и экспериментальная оценка эффекта каждого
механизма относительно базового варианта.

Для достижения поставленной цели необходимо:
\begin{enumerate}
  \item Провести обзор существующих решений для управления федерацией
        кластеров, наблюдаемости, аварийного восстановления и контроля
        политик размещения.
  \item Развернуть экспериментальный стенд из управляющего кластера и двух
        кластеров федерации с тестовым сервисом и генератором нагрузки.
  \item Построить слой наблюдаемости федерации: агрегацию метрик кластеров,
        показатели SLI/SLO и сквозную трассировку запросов.
  \item Разработать каталог кластеров, хранящий их атрибуты, текущее
        состояние и журнал принятых решений.
  \item Реализовать двухуровневые политики размещения: фильтрацию на уровне
        федерации и контроль внутри каждого кластера.
  \item Разработать планировщик уровня федерации, учитывающий текущие
        метрики кластеров, и наивный планировщик для сравнения.
  \item Разработать контроллер аварийного восстановления, автоматически
        выполняющий цикл «обнаружение~--- переключение трафика~---
        восстановление».
  \item Провести эксперименты по трём сценариям отказов и сравнить каждый
        механизм с вариантом без него.
  \item Проанализировать результаты: проверить гипотезы, оценить
        чувствительность к параметрам и достоверность измерений.
\end{enumerate}
